\scanpage{194}\mspage{143}
(et il reste à vérifier qu'on a ainsi un syst. \uline{transitif} d'isom. canoniques dans $\mathrm{HOT}_W(\Gamma)$, entre les quatre objets envisagés, \deleted{\unsure{)}} ce qui doit être essentiellement \unsure{sorital}\ldots).

\uline{Remarques} 1.) Ce qu'on vient de dire dans le cas des diagrammes en $\mathrm{Cat}$ de type $\Phi^\circ=\Psi$, \inserted{\unsure{a|à} l'air de} garder un sens pour \unsure{les|des} diagrammes de type quelconque

\[
X_*:S^\circ\longrightarrow\mathrm{Cat},\qquad s\longmapsto X_s .
\]

($S$ un objet de $\mathrm{Cat}$). On peut l'intégrer de façon lisse, $\int_{\deleted{\unsure{\ell}}}X_*\underset{\unsure{\mathrm{Df}|{D}}}{=}\int_{\mathrm{lis}}X_*$, obtenant une catégorie fibrée (scindée) sur $S$, ou de façon propre, $\int_{\mathrm{prop}}X_*$, obtenant une catégorie cofibrée scindée sur $S^\circ$. On peut également considérer $X_*^\circ$, composé de $X_*$ avec l'autom. $Y\longmapsto Y^\circ$ de $\mathrm{Cat}$, et l'intégrer de façon soit propre, soit lisse. Dans tous ces cas, on trouve un diagramme de type $\widehat S^\circ$, où $\widehat S$ est déduit de $S$ en \unsure{appendant} un élément \deleted{initial} \deleted{\unsure{décrit}} \deleted{\unsure{caractérisé par}} \inserted{qui} \deleted{une flèche \unsure{unique} dans tous les autres (sans} s'envoie par une flèche ``canonique'' \inserted{(\unsure{choisie|donnée|libre?})} dans chacun des autres (\uline{sans} relations de commutation !). Un
% archive sheet number: 194
